% fig_mechanism.tex -- Gaussian scale-mixture mechanism.
%
% One deliberately simple three-point p(t) is used everywhere in the figure:
%   t = (2.2, 5.0, 9.5),  w = (0.42, 0.40, 0.18),  <t> = 4.634,  a = 1.
% Panel (a): a porous slab.  Pores sit ON each ray, so the ray visibly threads
% them (thin in a pore/void, thick in solid).  The three rays cross a different
% NUMBER of pores (6 / 4 / 2) -- the honest picture of "different transverse
% position -> different solid path length t".  At exit each ray carries a
% sideways Gaussian of width sqrt(a t): short path -> narrow, long path -> wide.
% The bars below print the same discrete (t_i, w_i).  Panel (b) is exactly the
% corresponding weighted mixture and its equal-variance Gaussian reference.
%
% Layout: the slab (panel a) and the plot box (panel b) share top y=58 and
% bottom y=18; a clear column of white space separates the panels; every pore
% keeps a margin from all four slab edges (no clipping).
%
% Build from figs/src:
%   pdflatex -interaction=nonstopmode -halt-on-error -output-directory=../ fig_mechanism.tex
\documentclass[tikz]{standalone}
\usepackage[T1]{fontenc}
\usepackage{newtxtext,newtxmath}
\usepackage{pgfplots}
\pgfplotsset{compat=1.18}
\usetikzlibrary{arrows.meta,calc}

\definecolor{shortc}{HTML}{7EA4C9}
\definecolor{mediumc}{HTML}{4A76A6}
\definecolor{longc}{HTML}{243F5C}
\definecolor{hero}{HTML}{1A1A1A}
\definecolor{mute}{HTML}{666C73}
\definecolor{slabfill}{HTML}{E7EAEF}
\definecolor{poreedge}{HTML}{AEB7C2}
\definecolor{trackbg}{HTML}{E4E8ED}

% The illustrative distribution: these macros drive both panels.
\def\tA{2.2}
\def\tB{5.0}
\def\tC{9.5}
\def\wA{0.42}
\def\wB{0.40}
\def\wC{0.18}
\pgfmathsetmacro{\tMean}{\wA*\tA + \wB*\tB + \wC*\tC}

% The standalone is included at essentially native width in JINST.  Keep every
% printed label at or above 8.5 pt at final size.
\newcommand{\figlab}[1]{{\fontsize{9.0}{10.6}\selectfont #1}}
\newcommand{\figsmall}[1]{{\fontsize{8.6}{10.1}\selectfont #1}}

\begin{document}
\begin{tikzpicture}[x=1mm,y=1mm]

% ======================== (a) porous medium + conditional widths ===========
\node[anchor=base west,text=hero] at (0,60.3) {\fontsize{10.5}{12}\selectfont\textbf{(a)}};
% side labels: entry on the left, exit on the right, centred on the slab band.
\node[text=hero,rotate=90,anchor=center] at (0.5,38) {\figlab{incident protons}};
\node[text=hero,rotate=-90,anchor=center] at (65.5,38) {\figlab{conditional kink}};

% -- the two-phase slab.  Every pore keeps a margin from all four edges.
\draw[hero,line width=0.30mm,fill=slabfill] (10,18) rectangle (54,58);
% background texture (between the rays): the medium is porous throughout.
\foreach \cx in {14,23,32,41,50}
  \fill[white,draw=poreedge,line width=0.18mm] (\cx,55) circle[radius=2.2];
\foreach \cx in {18,27,36,45}
  \fill[white,draw=poreedge,line width=0.18mm] (\cx,44) circle[radius=2.2];
\foreach \cx in {14,22,30,38,46}
  \fill[white,draw=poreedge,line width=0.18mm] (\cx,32) circle[radius=2.2];
\foreach \cx in {18,27,36,45}
  \fill[white,draw=poreedge,line width=0.18mm] (\cx,21) circle[radius=2.2];
% on-ray pores: 6 on the short ray, 4 on the medium, 2 on the long.
\foreach \cx in {15,22,29,36,43,50}
  \fill[white,draw=poreedge,line width=0.18mm] (\cx,50) circle[radius=2.2];
\foreach \cx in {17,28,39,49}
  \fill[white,draw=poreedge,line width=0.18mm] (\cx,38) circle[radius=2.2];
\foreach \cx in {28,46}
  \fill[white,draw=poreedge,line width=0.18mm] (\cx,26) circle[radius=2.2];

% -- three rays.  Each ray is thin in its own colour where it is inside a pore
% (void) and thick where it traverses SOLID.  Incident stub at left.
% Short ray, y=50: threads 6 pores -> mostly void -> short solid path.
\draw[shortc,line width=0.28mm] (10,50)--(54,50);
\foreach \xa/\xb in {10/12.8,17.2/19.8,24.2/26.8,31.2/33.8,38.2/40.8,45.2/47.8,52.2/54}
  \draw[shortc,line width=0.9mm,line cap=butt] (\xa,50)--(\xb,50);
\draw[shortc,line width=0.42mm] (3,50)--(10,50);
% Medium ray, y=38: threads 4 pores -> intermediate solid path.
\draw[mediumc,line width=0.28mm] (10,38)--(54,38);
\foreach \xa/\xb in {10/14.8,19.2/25.8,30.2/36.8,41.2/46.8,51.2/54}
  \draw[mediumc,line width=0.9mm,line cap=butt] (\xa,38)--(\xb,38);
\draw[mediumc,line width=0.42mm] (3,38)--(10,38);
% Long ray, y=26: threads 2 pores -> nearly all solid path.
\draw[longc,line width=0.28mm] (10,26)--(54,26);
\foreach \xa/\xb in {10/25.8,30.2/43.8,48.2/54}
  \draw[longc,line width=0.9mm,line cap=butt] (\xa,26)--(\xb,26);
\draw[longc,line width=0.42mm] (3,26)--(10,26);

% -- conditional kink: a sideways Gaussian at each exit, width sqrt(a t_i);
% the vertical spread grows as sqrt(t_i).  All spreads stay within the slab band.
\def\amp{5.6}
\fill[shortc,opacity=0.14]
  plot[domain=-3.0:3.0,samples=41,smooth]
  ({54+\amp*exp(-(\x)*(\x)/(2*1.483*1.483))},{50+\x}) -- (54,53) -- (54,47) -- cycle;
\draw[shortc,line width=0.30mm]
  plot[domain=-3.0:3.0,samples=41,smooth]
  ({54+\amp*exp(-(\x)*(\x)/(2*1.483*1.483))},{50+\x});
\draw[shortc,line width=0.44mm,-{Stealth[length=1.8mm,width=1.5mm]}] (54,50)--(61.2,50);

\fill[mediumc,opacity=0.14]
  plot[domain=-4.5:4.5,samples=51,smooth]
  ({54+\amp*exp(-(\x)*(\x)/(2*2.236*2.236))},{38+\x}) -- (54,42.5) -- (54,33.5) -- cycle;
\draw[mediumc,line width=0.30mm]
  plot[domain=-4.5:4.5,samples=51,smooth]
  ({54+\amp*exp(-(\x)*(\x)/(2*2.236*2.236))},{38+\x});
\draw[mediumc,line width=0.44mm,-{Stealth[length=1.8mm,width=1.5mm]}] (54,38)--(61.2,38);

\fill[longc,opacity=0.14]
  plot[domain=-6.0:6.0,samples=61,smooth]
  ({54+\amp*exp(-(\x)*(\x)/(2*3.082*3.082))},{26+\x}) -- (54,32) -- (54,20) -- cycle;
\draw[longc,line width=0.30mm]
  plot[domain=-6.0:6.0,samples=61,smooth]
  ({54+\amp*exp(-(\x)*(\x)/(2*3.082*3.082))},{26+\x});
\draw[longc,line width=0.44mm,-{Stealth[length=1.8mm,width=1.5mm]}] (54,26)--(61.2,26);

% -- path-length bars: the bar length is t_i; the label prints its weight w_i.
\node[anchor=east,text=hero,align=right] at (8.5,6.5) {\figsmall{solid path}\\[-0.4ex]\figsmall{$t_i$}};
\foreach \yy in {11,6.5,2}
  \draw[trackbg,line width=1.4mm,line cap=round] (10,\yy)--(57.7,\yy);
\draw[shortc,line width=1.4mm,line cap=round]  (10,11)--({10+5*\tA},11);
\draw[mediumc,line width=1.4mm,line cap=round] (10,6.5)--({10+5*\tB},6.5);
\draw[longc,line width=1.4mm,line cap=round]   (10,2)--({10+5*\tC},2);
\node[anchor=west,text=shortc!75!black]  at (60,11)  {\figsmall{$t_1{=}2.2,\ w_1{=}0.42$}};
\node[anchor=west,text=mediumc!85!black] at (60,6.5) {\figsmall{$t_2{=}5.0,\ w_2{=}0.40$}};
\node[anchor=west,text=longc]            at (60,2)   {\figsmall{$t_3{=}9.5,\ w_3{=}0.18$}};
\draw[hero,line width=0.22mm,dash pattern=on 1.2mm off 0.9mm]
  ({10+5*\tMean},0.4)--({10+5*\tMean},12.6);
\node[anchor=north,text=hero] at ({10+5*\tMean},-0.6) {\figsmall{$\langle t\rangle=4.63=fL$}};

% ======================= (b) corresponding marginal =======================
% Plot box aligned to the slab: south-west (90,18), 62 x 40 mm -> top y=58.
% A clear white column separates it from panel (a)'s exit annotation.
\node[anchor=base west,text=hero] at (85,60.3) {\fontsize{10.5}{12}\selectfont\textbf{(b)}};

\begin{axis}[
  name=mainb,at={(90mm,18mm)},anchor=south west,
  scale only axis,width=62mm,height=40mm,
  xmin=-10.4,xmax=9.0,ymin=0,ymax=0.252,
  axis x line*=bottom,axis y line*=left,
  line width=0.26mm,axis line style={hero},
  xtick={-5,0,5},ytick=\empty,
  tick style={hero,line width=0.25mm},
  tick label style={font=\fontsize{8.7}{10}\selectfont},
  xlabel={\fontsize{9.2}{10.8}\selectfont projected kink angle $\theta$},
  ylabel={\fontsize{9.2}{10.8}\selectfont density},
  xlabel style={yshift=0.4mm},ylabel style={yshift=-1.2mm},
  clip=false]
  % Weighted components: each thin curve integrates to its printed w_i.
  \addplot[shortc,line width=0.30mm,domain=-9:9,samples=181,smooth]
    {\wA/sqrt(2*pi*\tA)*exp(-x^2/(2*\tA))};
  \addplot[mediumc,line width=0.30mm,domain=-9:9,samples=181,smooth]
    {\wB/sqrt(2*pi*\tB)*exp(-x^2/(2*\tB))};
  \addplot[longc,line width=0.30mm,domain=-9:9,samples=181,smooth]
    {\wC/sqrt(2*pi*\tC)*exp(-x^2/(2*\tC))};

  % Gaussian with identical variance a<t> (a=1 in the illustration).
  \addplot[mute,line width=0.34mm,dash pattern=on 1.5mm off 1.0mm,
           domain=-9:9,samples=181,smooth]
    {1/sqrt(2*pi*\tMean)*exp(-x^2/(2*\tMean))};

  % Exact weighted scale mixture defined above.
  \addplot[hero,line width=0.58mm,domain=-9:9,samples=201,smooth]
    {\wA/sqrt(2*pi*\tA)*exp(-x^2/(2*\tA))
     +\wB/sqrt(2*pi*\tB)*exp(-x^2/(2*\tB))
     +\wC/sqrt(2*pi*\tC)*exp(-x^2/(2*\tC))};

  % Frameless legend (no leader lines).  Short labels keep clear of the inset.
  \draw[hero,line width=0.58mm] (axis cs:-10.05,0.243)--(axis cs:-8.35,0.243);
  \node[anchor=west,text=hero,font=\fontsize{8.6}{10}\selectfont]
    at (axis cs:-8.15,0.243) {scale mixture};
  \draw[mute,line width=0.34mm,dash pattern=on 1.5mm off 1.0mm]
    (axis cs:-10.05,0.221)--(axis cs:-8.35,0.221);
  \node[anchor=west,text=mute,font=\fontsize{8.6}{10}\selectfont]
    at (axis cs:-8.15,0.221) {equal-variance};
  \draw[mediumc,line width=0.30mm] (axis cs:-10.05,0.199)--(axis cs:-8.35,0.199);
  \node[anchor=west,text=mute,font=\fontsize{8.6}{10}\selectfont]
    at (axis cs:-8.15,0.199) {components};
\end{axis}

% Log-density inset: same two curves; clear in the top-right corner.
\begin{axis}[
  at={($(mainb.north east)+(-0.6mm,-0.6mm)$)},anchor=north east,
  width=28mm,height=24mm,
  xmin=0,xmax=8.5,ymin=1e-4,ymax=0.30,ymode=log,
  axis background/.style={fill=white},
  axis x line*=bottom,axis y line*=left,
  line width=0.23mm,axis line style={hero},
  xtick={0,4,8},ytick={1e-1,1e-3},
  tick style={hero,line width=0.22mm},
  tick label style={font=\fontsize{8.3}{9.5}\selectfont,fill=white,inner sep=0.5pt},
  title={\fontsize{8.5}{9.8}\selectfont heavy tails},
  title style={yshift=-1.3mm}]
  \addplot[mute,line width=0.30mm,dash pattern=on 1.3mm off 0.9mm,
           domain=0:8.5,samples=121,smooth]
    {1/sqrt(2*pi*\tMean)*exp(-x^2/(2*\tMean))};
  \addplot[hero,line width=0.47mm,domain=0:8.5,samples=151,smooth]
    {\wA/sqrt(2*pi*\tA)*exp(-x^2/(2*\tA))
     +\wB/sqrt(2*pi*\tB)*exp(-x^2/(2*\tB))
     +\wC/sqrt(2*pi*\tC)*exp(-x^2/(2*\tC))};
\end{axis}

\end{tikzpicture}
\end{document}
